\begin{exercise}[Half risky, half risk-free strategy]
\emph{Restatement.}  In a one-risky-asset market, an investor always
rebalances so that half of the portfolio value is invested in the risky asset
and half in the risk-free asset.  Find the recursions for the holdings and for
the portfolio value.
\end{exercise}

\begin{solution}
Write
\[
  \theta_{t_k}=(\theta^0_{t_k},\theta^1_{t_k})
\]
for the holdings chosen at \(t_{k-1}\) and held over the interval
\((t_{k-1},t_k]\), as in \polyref{Def.~1.1}.  At time \(t_{k-1}\), before the
new holdings are chosen, the wealth is
\[
  V_{t_{k-1}}=\inner{\theta_{t_{k-1}}}{S_{t_{k-1}}}
  =\theta^0_{t_{k-1}}S^0_{t_{k-1}}
   +\theta^1_{t_{k-1}}S^1_{t_{k-1}}.
\]
The rule says that after rebalancing at the old prices \(S_{t_{k-1}}\),
\[
  \theta^0_{t_k}S^0_{t_{k-1}}=\frac12 V_{t_{k-1}},
  \qquad
  \theta^1_{t_k}S^1_{t_{k-1}}=\frac12 V_{t_{k-1}}.
\]
Therefore
\[
  \boxed{
  \theta^0_{t_k}
    =\frac{\inner{\theta_{t_{k-1}}}{S_{t_{k-1}}}}
           {2S^0_{t_{k-1}}},
  \qquad
  \theta^1_{t_k}
    =\frac{\inner{\theta_{t_{k-1}}}{S_{t_{k-1}}}}
           {2S^1_{t_{k-1}}}.
  }
\]
This is exactly the recurrence derived in \polyref{Chap.~4, Ex.~1.5}.

The value at the next date is obtained by keeping the holdings fixed while
prices move:
\[
\begin{aligned}
  V_{t_k}
    &=\theta^0_{t_k}S^0_{t_k}+\theta^1_{t_k}S^1_{t_k} \\
    &=\frac{V_{t_{k-1}}}{2S^0_{t_{k-1}}}S^0_{t_k}
      +\frac{V_{t_{k-1}}}{2S^1_{t_{k-1}}}S^1_{t_k}.
\end{aligned}
\]
Hence the wealth recursion is
\[
  \boxed{
  V_{t_k}
  =
  V_{t_{k-1}}
  \left(
    \frac12\frac{S^0_{t_k}}{S^0_{t_{k-1}}}
    +
    \frac12\frac{S^1_{t_k}}{S^1_{t_{k-1}}}
  \right).
  }
\]
Equivalently, if \(S^0_{t_k}/S^0_{t_{k-1}}=e^{r(t_k-t_{k-1})}\), then
\[
  V_{t_k}
  =
  V_{t_{k-1}}
  \left(
    \frac12 e^{r(t_k-t_{k-1})}
    +
    \frac12\frac{S^1_{t_k}}{S^1_{t_{k-1}}}
  \right).
\]
\end{solution}
